\documentclass[12pt,a4paper]{article}
\usepackage[utf8]{inputenc}
\usepackage[margin=1in]{geometry}
\usepackage{graphicx}
\usepackage{enumitem}
\usepackage{hyperref}
\usepackage{xcolor}
\usepackage{booktabs}
\usepackage{fancyhdr}
\usepackage{titlesec}
\usepackage{listings}
\usepackage{float}
\usepackage{caption}

\definecolor{paleoblue}{RGB}{74,158,255}
\definecolor{paleogreen}{RGB}{76,175,80}
\definecolor{paleored}{RGB}{200,50,50}
\definecolor{darkgray}{RGB}{80,80,80}
\definecolor{codebg}{RGB}{245,245,245}

\hypersetup{
  colorlinks=true,
  linkcolor=paleoblue,
  urlcolor=paleoblue,
}

\pagestyle{fancy}
\fancyhf{}
\fancyhead[L]{\small PaleoVox GUI Documentation}
\fancyhead[R]{\small v1.0.6}
\fancyfoot[C]{\thepage}

\titleformat{\section}{\Large\bfseries}{\thesection}{1em}{}[\color{paleoblue}\titlerule]
\titleformat{\subsection}{\large\bfseries}{\thesubsection}{1em}{}

\lstset{
  basicstyle=\ttfamily\small,
  backgroundcolor=\color{codebg},
  frame=single,
  rulecolor=\color{gray!40},
  breaklines=true,
  showstringspaces=false,
  columns=flexible,
}

\begin{document}

\begin{titlepage}
\centering
\vspace*{3cm}

{\Huge\bfseries PaleoVox}\\[0.5cm]
{\Large Graphical User Interface}\\[0.5cm]
{\Large Documentation}\\[1.5cm]

{\large Version 1.0.6}\\[2cm]

\begin{tabular}{rl}
\textbf{Date:} & \today \\
\textbf{Authors:} & Alan Gabriel Amaro Colin \\
                 & Ángel Ángeles Cortés \\
                 & Jair Israel Barrientos Lara \\
                 & Jesus Alberto Díaz Cruz \\
\textbf{Institution:} & Facultad de Ciencias \& Instituto de Geología, UNAM \\
\end{tabular}

\vfill
\end{titlepage}

\tableofcontents
\newpage

% =============================================================================
\section{Introduction}
% =============================================================================

\subsection{What is PaleoVox?}

PaleoVox is a Python-based application for 3D fossil data augmentation using voxel representations. It provides a graphical interface for:

\begin{itemize}
  \item Converting 3D meshes to voxel grids
  \item Applying morphological operations (dilation/erosion)
  \item Simulating taphonomic damage (deformation, erosion, rotation, fracture)
  \item Reconstructing high-quality meshes from voxels
  \item Visualizing and comparing original vs. processed data
  \item Dimensionality reduction analysis via t-SNE
  \item 2D perspective projections
\end{itemize}

\subsection{Supported File Formats}

\begin{table}[H]
\centering
\begin{tabular}{llp{7cm}}
\toprule
\textbf{Extension} & \textbf{Type} & \textbf{Description} \\
\midrule
\texttt{.ply} & Mesh & Polygon File Format with vertex colors \& normals \\
\texttt{.obj} & Mesh & Wavefront OBJ, widely compatible \\
\texttt{.npy}  & Voxels & NumPy binary array format for voxel grids \\
\bottomrule
\end{tabular}
\end{table}

\subsection{System Requirements}
\begin{itemize}
  \item Python 3.8 or later
  \item PyQt5, Open3D, NumPy, SciPy, Matplotlib, scikit-learn
  \item X11 display server (Linux)
  \item Minimum screen resolution: 1100$\times$700
\end{itemize}

\subsection{Launching PaleoVox}
To start the application, run from the project directory:
\begin{lstlisting}[language=bash]
python3 paleovox_gui.py
\end{lstlisting}

\newpage

% =============================================================================
\section{Interface Overview}
% =============================================================================

The PaleoVox window is divided into two main panels separated by a movable splitter:

\begin{description}
  \item[\textbf{Left Panel} (\~{}350px)] Contains the logo, file Drop Zone, Browse button, File Info panel, Reset button, and About link.
  \item[\textbf{Right Panel} (\~{}700px)] Contains five tabbed workspaces for the core workflow.
\end{description}

\begin{table}[H]
\centering
\begin{tabular}{lp{10cm}}
\toprule
\textbf{Tab} & \textbf{Purpose} \\
\midrule
Pipeline \& Augmentation & Core workflow: mesh loading, voxel conversion, dilation, and four augmentation operations \\
Reconstruction \& Comparison & Mesh reconstruction from voxels, visual comparison of meshes and voxels \\
View \& Save & 3D visualization and file export \\
t-SNE Analysis & Dimensionality reduction to compare original vs. deformed voxel distributions \\
2D Perspectives & 2D projection scatter plots of voxel data \\
\bottomrule
\end{tabular}
\end{table}

% =============================================================================
\section{Left Panel --- File Management}
% =============================================================================

\subsection{Drop Zone}

The Drop Zone is a large dashed-bordered area that accepts 3D files via drag-and-drop or click-to-browse. It has three visual states:

\begin{description}
  \item[\textbf{Idle} (gray)] No file loaded. Displays: ``Drop .ply, .obj or .npy file here''.
  \item[\textbf{Hover} (blue)] A supported file is being dragged over the zone.
  \item[\textbf{Loaded} (green)] A file has been successfully loaded. Shows the filename.
\end{description}

\subsection{Browse Button}

Clicking ``Browse...'' opens a file dialog filtered to supported formats. The same dialog is opened when clicking directly on the Drop Zone text.

\subsection{File Info Panel}

Once a file is loaded, the File Info panel displays:

\begin{itemize}
  \item \textbf{Path:} Full file path (with home directory abbreviated as \texttt{\~{}})
  \item \textbf{Vertices:} Number of mesh vertices (``---'' for \texttt{.npy} files)
  \item \textbf{Triangles:} Number of mesh triangles (``---'' for \texttt{.npy} files)
  \item \textbf{Voxel shape:} Dimensions of the voxel grid (e.g., \texttt{(256, 256, 256)})
  \item \textbf{Occupied voxels:} Count of non-zero voxels
\end{itemize}

\subsection{Reset All}

Clears all data and returns the application to its initial state. All visualizations, augmentations, and reconstructions are discarded.

\subsection{More about PaleoVox}

Opens an About dialog showing the application description and creator credits.

% =============================================================================
\section{Tab 1: Pipeline \& Augmentation}
% =============================================================================

This tab contains two grouped sections: the Pipeline and the Augmentations.

\subsection{Pipeline Section}

\begin{table}[H]
\centering
\begin{tabular}{llp{8cm}}
\toprule
\textbf{Control} & \textbf{Type} & \textbf{Description} \\
\midrule
Load Mesh & Button & Re-loads the currently loaded file \\
npoints & SpinBox (100--200000) & Number of points sampled from mesh surface \\
dim & SpinBox (32--512) & Voxel grid resolution (N$\times$N$\times$N) \\
Mesh $\rightarrow$ Voxels & Button & Converts the current mesh to a voxel grid \\
Iterations & SpinBox (1--10) & Dilation iterations for morphological closing \\
Dilate Voxels & Button & Applies binary dilation to the voxel array \\
Voxels $\rightarrow$ Mesh & Button & Reconstructs a mesh from the current voxels \\
\bottomrule
\end{tabular}
\end{table}

\subsubsection{Workflow}
The typical pipeline order is:
\begin{enumerate}
  \item \textbf{Load Mesh} --- Drag/set a file. Mesh info appears in the left panel.
  \item \textbf{Mesh $\rightarrow$ Voxels} --- Convert the mesh to a 3D binary voxel grid. Set \texttt{npoints} (sampling density) and \texttt{dim} (resolution) before converting. Higher \texttt{npoints} captures more surface detail; higher \texttt{dim} creates finer grids at the cost of memory and speed.
  \item \textbf{Dilate Voxels} --- (Optional) Fill interior gaps and close surface holes. Higher iterations produce more aggressive filling.
  \item \textbf{Voxels $\rightarrow$ Mesh} --- Reconstruct a triangle mesh from voxels using Poisson surface reconstruction with post-processing.
\end{enumerate}

\subsection{Augmentations Section}

All augmentations operate on the \textbf{current voxel grid}. The original voxel grid is preserved separately for comparison.

\subsubsection{Deformation}

Simulates sediment compaction along a chosen axis.

\begin{itemize}
  \item \textbf{Axis:} X, Y, or Z (default: Z for vertical compaction)
  \item \textbf{Factor:} 0.10--0.99 (default: 0.85). Lower values $\rightarrow$ stronger compression. A factor of 1.0 means no change.
\end{itemize}

\subsubsection{Erosion}

Simulates partial degradation by removing voxels beyond a random cut point along a chosen axis.

\begin{itemize}
  \item \textbf{Axis:} X, Y, or Z
  \item \textbf{Increment:} 0.01--1.0 (default: 0.5). Controls minimum preservation. Higher values keep more material.
\end{itemize}

\subsubsection{Rotation}

Rigid 3D rotation of the voxel grid about its geometric center. Nearest-neighbor interpolation preserves binary values.

\begin{itemize}
  \item \textbf{X/Y/Z (degrees):} $-180$ to $180$. Positive values follow the right-hand rule.
\end{itemize}

\subsubsection{Fracture}

Stochastic fracture propagation through the voxel volume. Simulates cracks by propagating from a starting point through occupied voxels.

\begin{itemize}
  \item \textbf{Max pos:} 1--50 (default: 10). Controls randomness. Lower values $\rightarrow$ more directed fractures; higher values $\rightarrow$ more wandering.
  \item \textbf{Return both:} If checked, also extracts the fracture pattern as a separate grid. Status bar reports pattern voxel count.
\end{itemize}

\subsubsection{Save Deformed Voxels}

Exports the current voxel grid as a \texttt{.npy} file. Useful for saving intermediate states or results for later use.

% =============================================================================
\section{Tab 2: Reconstruction \& Comparison}
% =============================================================================

\subsection{Reconstruction \& Comparison Group}

\begin{table}[H]
\centering
\begin{tabular}{llp{8cm}}
\toprule
\textbf{Control} & \textbf{Type} & \textbf{Description} \\
\midrule
Reconstructed: & Label & Displays vertex/triangle count of the reconstructed mesh \\
Reconstruct from Voxels & Button & Runs Poisson surface reconstruction from the current voxels \\
Compare Original vs Reconstructed & Button & Opens a 3D viewer with the original and reconstructed meshes \\
Show: & ComboBox & Filters the comparison: Both, Original Only, or Reconstructed Only \\
Save Reconstructed Mesh & Button & Exports the reconstructed mesh as \texttt{.ply} or \texttt{.obj} \\
\bottomrule
\end{tabular}
\end{table}

\subsubsection{Reconstruction Details}
Reconstruction uses high-quality Poisson surface reconstruction with:
\begin{itemize}
  \item Morphological solidification (closing + hole filling)
  \item Positional jitter to break grid regularity
  \item Statistical outlier removal (SOR)
  \item Density-based filtering to remove low-confidence geometry
  \item Connected component filtering to remove floating artifacts
  \item Quadric edge collapse decimation (if needed, capped at 100K triangles)
  \item Taubin smoothing for surface quality
\end{itemize}

If the original mesh scale information is available (from a mesh-loaded workflow), the reconstructed mesh is scaled and translated to match the original bounds.

\subsection{Voxel Comparison Group}

\begin{table}[H]
\centering
\begin{tabular}{llp{7cm}}
\toprule
\textbf{Control} & \textbf{Type} & \textbf{Description} \\
\midrule
Original: --- Current: --- & Label & Voxel counts for original and current grids \\
Show: & ComboBox & Visualization filter: Both, Original Only, Current Only \\
Compare Voxels & Button & Opens a 3D viewer with overlaid voxel grids \\
\bottomrule
\end{tabular}
\end{table}

The voxel comparison visualizes original (blue) and current (red) voxel grids as 3D point clouds. This directly shows which regions were added or removed by augmentations.

% =============================================================================
\section{Tab 3: View \& Save}
% =============================================================================

\subsection{Color Selector}

A global color picker that sets the display color for both mesh and voxel 3D viewers. Available colors: Blue, Red, Green, Orange, Purple, Cyan, Yellow, White, Gray (default: Blue).

\subsection{View Buttons}

\begin{itemize}
  \item \textbf{View Mesh} --- Opens an Open3D interactive 3D viewer showing the current or reconstructed mesh. Supports rotation, pan, zoom.
  \item \textbf{View Voxels} --- Opens an Open3D viewer showing the current voxel grid as a voxelized point cloud.
\end{itemize}

Both viewers open in separate windows and run in background threads, so the main GUI remains responsive.

\subsection{Save Buttons}

\begin{itemize}
  \item \textbf{Save Mesh} --- Exports the current mesh as \texttt{.ply} or \texttt{.obj}.
  \item \textbf{Save Voxel} --- Exports the current voxel grid as a \texttt{.npy} NumPy file.
\end{itemize}

Default filenames are derived from the source file name (e.g., \texttt{fossil\_processed.ply}).

% =============================================================================
\section{Tab 4: t-SNE Analysis}
% =============================================================================

t-Distributed Stochastic Neighbor Embedding (t-SNE) projects the 3D coordinates of occupied voxels into a 2D space for visualization. It reveals clustering patterns and distribution differences between original and deformed voxel grids.

\subsection{Parameters}

\begin{table}[H]
\centering
\begin{tabular}{lllp{6cm}}
\toprule
\textbf{Parameter} & \textbf{Range} & \textbf{Default} & \textbf{Description} \\
\midrule
Perplexity & 5--500 & 100 & Balance between local/global structure \\
Seed & 0--9999 & 42 & Random seed for reproducibility \\
Percentage & 0.01--1.00 & 0.50 & Fraction of occupied voxels to sample \\
Point size & 0.1--10.0 & 1.0 & Scatter marker size \\
\bottomrule
\end{tabular}
\end{table}

\subsection{Output}

Clicking ``Generate t-SNE'' produces a dialog with two side-by-side scatter plots:
\begin{itemize}
  \item \textbf{Left:} Original voxels in blue
  \item \textbf{Right:} Current (deformed) voxels in red
\end{itemize}

The embeddings are computed in a shared space by concatenating both point sets, fitting t-SNE once, then splitting the result. This ensures direct comparability.

\textbf{Note:} Both the original and current voxel grids must contain occupied voxels. An error dialog appears if either is empty.

\subsection{Saving}

The ``Save Image...'' button in the dialog exports the figure as PNG or JPEG. The temp file is automatically cleaned up when the dialog is closed.

% =============================================================================
\section{Tab 5: 2D Perspectives}
% =============================================================================

This tab projects all occupied voxels onto a 2D plane, creating an X-ray view of the voxel distribution. It is useful for quickly inspecting shape, density, and orientation from different angles.

\subsection{Single Perspective View}

\begin{table}[H]
\centering
\begin{tabular}{lllp{7cm}}
\toprule
\textbf{Control} & \textbf{Type} & \textbf{Default} & \textbf{Description} \\
\midrule
Axis & ComboBox & XY & Projection plane (XY, XZ, or YZ) \\
Color & ComboBox & Blue & Scatter point color \\
Marker & ComboBox & o & Marker style (o, ., \^{}, s, *, +, x) \\
Size & SpinBox & 1.0 & Marker size (0.1--10.0) \\
\bottomrule
\end{tabular}
\end{table}

Click ``Generate 2D Perspective'' to produce a scatter plot projection of the current voxel grid onto the selected plane. For example, XY projection shows the object as seen from above, collapsing all Z-depth information into a single view.

\subsection{Comparison View (Original vs Current)}

Overlays both the original and current voxel projections on the same axes, with independently customizable colors, markers, and sizes for each.

\begin{table}[H]
\centering
\begin{tabular}{llp{7cm}}
\toprule
\textbf{Control} & \textbf{Default} & \textbf{Description} \\
\midrule
Orig color / Curr color & Blue / Red & Colors for original and current voxels \\
Marker (each) & o / \^{} & Marker styles for each set \\
Size (each) & 1.0 / 1.0 & Marker sizes for each set \\
\bottomrule
\end{tabular}
\end{table}

Click ``Generate Comparison'' to produce an overlaid scatter plot. A legend distinguishes the two datasets. This is ideal for evaluating the effect of augmentations.

\subsection{Saving}

Both dialogs include a ``Save Image...'' button to export the plot as PNG or JPEG.

% =============================================================================
\section{Working with \texttt{.npy} Files Directly}
% =============================================================================

PaleoVox supports loading pre-computed voxel grids directly from \texttt{.npy} files, skipping the mesh conversion step.

\subsection{Loading}
Drag a \texttt{.npy} file onto the Drop Zone or use Browse. The voxel grid is loaded immediately.

\subsection{State After Loading}
\begin{itemize}
  \item \textbf{Voxels:} Set to the loaded array. Original voxels are also captured.
  \item \textbf{Mesh:} Not available. You must reconstruct from voxels first.
  \item \textbf{Augmentations:} All enabled immediately (they operate on voxels).
  \item \textbf{Scale info:} Not available. Reconstructed mesh uses unit scale.
\end{itemize}

\subsection{Enabling Mesh Features}
Click ``Reconstruct from Voxels'' or ``Voxels $\rightarrow$ Mesh'' to generate a mesh. After reconstruction, all mesh features (View Mesh, Save Mesh, mesh comparison) become available.

\subsection{Common Use Case}
\begin{enumerate}
  \item Load a \texttt{.npy} file saved from a previous PaleoVox session.
  \item Apply augmentations (deformation, erosion, etc.).
  \item Reconstruct a mesh.
  \item Save the reconstructed mesh as \texttt{.ply} or \texttt{.obj}.
  \item Save the augmented voxels as \texttt{.npy} for later use.
\end{enumerate}

% =============================================================================
\section{Button State Logic}
% =============================================================================

PaleoVox dynamically enables and disables buttons based on the current application state to prevent invalid operations. The logic follows the principle that each operation requires its inputs to be present.

\begin{table}[H]
\centering
\small
\begin{tabular}{lccccc}
\toprule
\textbf{Button} & \textbf{Start} & \textbf{Mesh} & \textbf{Voxels} & \textbf{Recon} & \textbf{.npy} \\
\midrule
Load Mesh                   & -- & Y & Y & Y & Y \\
Mesh $\rightarrow$ Voxels  & -- & Y & -- & -- & -- \\
Dilate Voxels               & -- & -- & Y & Y & Y \\
Voxels $\rightarrow$ Mesh  & -- & -- & Y & Y & Y \\
\addlinespace
All Augmentations           & -- & -- & Y & Y & Y \\
\midrule
View Mesh                   & -- & Y & -- & Y & -- \\
View Voxels                 & -- & -- & Y & Y & Y \\
Save Mesh                   & -- & Y & -- & Y & -- \\
Save Voxel                  & -- & -- & Y & Y & Y \\
\midrule
Reconstruct from Voxels     & -- & -- & Y & Y & Y \\
Compare Orig vs Recon       & -- & -- & -- & Y* & -- \\
Save Reconstructed          & -- & -- & -- & Y & -- \\
Compare Voxels              & -- & -- & Y** & Y** & Y** \\
Generate t-SNE              & -- & -- & Y** & Y** & Y** \\
\midrule
2D Single Perspective       & -- & -- & Y & Y & Y \\
2D Comparison               & -- & -- & Y** & Y** & Y** \\
\bottomrule
\multicolumn{6}{l}{\footnotesize Start = no file; Mesh = .ply/.obj loaded; Voxels = voxelized; Recon = mesh rebuilt; .npy = .npy loaded} \\
\multicolumn{6}{l}{\footnotesize Y = enabled; -- = disabled; * = requires original mesh; ** = requires original voxels} \\
\end{tabular}
\normalsize
\end{table}

% =============================================================================
\section{Error Handling}
% =============================================================================

PaleoVox provides user-friendly error handling throughout:

\begin{description}
  \item[\textbf{Load Errors}] If a file cannot be read (corrupted, wrong format), an error dialog appears with details. The application does not crash.
  \item[\textbf{Operation Errors}] If a backend operation fails (e.g., reconstruction of an empty voxel grid), an error dialog explains the issue.
  \item[\textbf{Save Errors}] If a file cannot be written (e.g., read-only directory), an error dialog shows the cause.
  \item[\textbf{Unsupported Files}] Drag-and-drop silently ignores unsupported file types. The file dialog filters prevent selection of unsupported formats.
  \item[\textbf{Empty Data}] Viewing an empty voxel grid shows a status message rather than opening a blank viewer.
\end{description}

% =============================================================================
\section{Tips \& Best Practices}
% =============================================================================

\subsection{Mesh Preparation}
\begin{itemize}
  \item Ensure meshes are clean (no degenerate triangles, non-manifold edges). Open3D handles minor issues but poorly formed meshes may produce unexpected results.
  \item Meshes should be watertight for best reconstruction quality.
\end{itemize}

\subsection{Voxel Parameters}
\begin{itemize}
  \item \textbf{npoints:} Increase for meshes with fine surface detail. Values above 100K may cause slower conversion.
  \item \textbf{dim:} Higher resolution (e.g., 512) captures finer detail but uses significantly more memory. 128 or 256 is suitable for most fossils.
\end{itemize}

\subsection{Augmentation Workflow}
\begin{itemize}
  \item Apply augmentations in any order. Each operates on the current voxel state.
  \item Use ``Reset All'' to discard all augmentations and start over.
  \item Save intermediate voxel states (\texttt{.npy}) if you want to revert to a specific point.
\end{itemize}

\subsection{Reconstruction Quality}
\begin{itemize}
  \item Reconstruction works best on solid (dilated) voxel grids. Use dilation before reconstruction to close gaps.
  \item For \texttt{.npy}-loaded voxels without scale info, the reconstructed mesh uses arbitrary units. Scale in your target application if needed.
\end{itemize}

\subsection{Visual Comparison}
\begin{itemize}
  \item Use the Voxel Comparison view to see exactly which regions were affected by augmentations.
  \item Use the Mesh Comparison view to evaluate reconstruction accuracy vs. the original.
  \item Use t-SNE to detect distribution-level changes that may not be visually obvious.
  \item Use 2D Perspectives for quick shape inspections from different angles.
\end{itemize}

% =============================================================================
\section{Keyboard Shortcuts \& Interactions}
% =============================================================================

\begin{table}[H]
\centering
\begin{tabular}{ll}
\toprule
\textbf{Action} & \textbf{Method} \\
\midrule
Load file & Drag \& drop or click Browse \\
Reset application & Click ``Reset All'' \\
3D viewer (rotate/pan/zoom) & Standard Open3D mouse controls \\
Close dialogs & Click ``Close'' or press Escape \\
\bottomrule
\end{tabular}
\end{table}

\subsection{Open3D Viewer Controls}
\begin{itemize}
  \item \textbf{Left mouse drag:} Rotate the view
  \item \textbf{Middle mouse drag / Scroll:} Zoom in/out
  \item \textbf{Right mouse drag / Ctrl + Left drag:} Pan
  \item \textbf{R:} Reset view to default
  \item \textbf{W:} Toggle wireframe mode
  \item \textbf{S:} Toggle shading mode
\end{itemize}

\end{document}
